\section*{Terminology and notation policy}

Literature in formal logic does not impose one universal notation.  Pretending
otherwise would be inaccurate.  This document instead fixes a registered
global convention.  A subsystem-specific departure is permitted only when its
scope and standing are declared before use, and no registered notation is used
for a second role within that scope.

Whenever this document introduces a new technical name, five conditions are
met at the same location:
\begin{enumerate}
  \item \emph{Necessity}: an existing term cannot express the required
        distinction without ambiguity;
  \item \emph{Definition}: an exact eliminable definition or an explicit axiom
        is supplied;
  \item \emph{Standing}: the grammatical and formal standing of the
        name is stated;
  \item \emph{Nonredundancy}: the distinction from previously defined terms is
        explained; and
  \item \emph{Use}: later definitions or theorems require the distinction.
\end{enumerate}
Until those conditions are met, the text describes the required object or
relation without assigning it a technical name.

Throughout this document, the word \emph{iff} abbreviates the
metalinguistic biconditional ``if and only if''; it is the registered
connective of definitions and metatheoretic equivalences.

The following distinctions are permanent:
\begin{itemize}
  \item A \emph{formula} is a grammatical object-language expression.  A
        \emph{sentence} is a formula with no free variable occurrence.
  \item A \emph{structure} supplies a domain and interpretations for the
        nonlogical symbols of a signature.  A \emph{model} is a structure that
        satisfies a specified set of sentences.
  \item A \emph{valuation} assigns truth values to propositional variables.  A
        \emph{variable assignment} assigns domain elements to first-order
        variables.
  \item \emph{Satisfaction} is the relation between a formula and the complete
        semantic data permitted by the active semantics.  \emph{Truth} is
        satisfaction relative to specified such data.  \emph{Validity} is
        satisfaction for every permitted tuple of semantic data.
        \emph{Semantic consequence} is preservation of satisfaction from a set
        of assumptions to a conclusion.  For a first-order sentence, assignment
        independence later permits the specialized phrase ``truth in a
        structure.''
  \item A \emph{derivation} is a finite rule-governed object in a specified
        calculus.  A \emph{proof} is a derivation with no undischarged
        assumptions other than the stated premises.
  \item \emph{Soundness} proceeds from derivability to semantic consequence.
        \emph{Completeness} proceeds from semantic consequence to derivability.
  \item \emph{Consistency} excludes derivability of a contradiction.
        A theory is \emph{deductively complete} iff every sentence or its
        negation is derivable from the theory.  A theory is
        \emph{semantically complete} iff it semantically entails every sentence
        or its negation.  For the fixed sound and complete first-order calculus,
        these two properties are equivalent; neither is consistency.
\end{itemize}

